\chapter*{ABSTRACT}
\addcontentsline{toc}{chapter}{ABSTRACT}

\begin{spacing}{1.5}
Phishing represents one of the most pervasive and rapidly evolving cyber threats in modern digital ecosystems. The proliferation of automated Phishing-as-a-Service (PaaS) platforms has substantially diminished the technical barrier required to execute distributed social engineering attacks, disproportionately endangering critical sectors including banking, financial technology, and enterprise communications. Traditional defense mechanisms, comprising static URL blacklists, domain reputation databases, and heuristic filters, exhibit critical limitations against polymorphic obfuscation, typosquatting, and fast-flux zero-day phishing campaigns, resulting in high false-negative rates and delayed threat mitigation.

While recent advancements in machine learning and deep neural networks have improved automated detection accuracy, their internal ``black-box'' decision-making characteristics impose severe impediments to adoption within enterprise Security Operations Centers (SOCs). Modern cybersecurity analysts require transparent, human-interpretable justifications before executing containment protocols such as domain blacklisting or network isolation. Furthermore, unimodal detection frameworks that analyze only lexical URL structures, raw HTML code, or rendered page visual appearances remain intrinsically vulnerable to evasion when an attacker obscures one specific vector while mimicking benign targets across others.

To address these core challenges, this seminar examines an advanced multimodal late-fusion phishing website detection architecture integrated with explainable artificial intelligence (XAI) technologies. The base system processes four complementary data modalities through specialized, independent branches: a CatBoost model for engineered URL features and domain metadata, a 1D Convolutional Neural Network (1D-CNN) for symbolic character-level sequence representations, a fine-tuned CodeBERT transformer for semantic and structural HTML source code parsing, and an EfficientNet-B7 deep convolutional backbone for screenshot-based visual layout and brand imitation analysis. The outputs of these specialized branches are synthesized through a trainable late-fusion meta-classifier optimized for the minority phishing class.

To ensure actionable operational interpretability, the architecture integrates a multi-tiered XAI suite comprising TreeSHAP for tabular URL features, 1D Integrated Gradients and attention matrix extraction for code and character sequences, and Grad-CAM for visual screenshot heatmaps. A local Large Language Model (LLM) aggregates these quantitative explanations into an analyst-oriented natural language incident report, augmented with threat intelligence context. To reinforce this research, three supporting investigations are reviewed, analyzing cross-channel behavior-based detection via EM-BERT and SPCA-based explainable LSTMs, dynamic arithmetic optimization for character-level URL detection in cloud environments (DAOA-DLPC), and comparative evaluations of recurrent and convolutional architectures (LSTM, BiLSTM, GRU, BiGRU).

Experimental evaluations demonstrate that the late-fusion multimodal paradigm achieves state-of-the-art performance, recording F1-scores of 0.989 on a 3.2 million URL enterprise corpus and 0.953 on the public MTLP benchmark dataset, while successfully identifying 92 out of 100 live zero-day phishing links, significantly outperforming Google Safe Browsing and OpenPhish. SOC analyst evaluations confirm high usability, scoring 3.97 out of 5.0, which further elevates to 4.20 with retrieval-augmented context. The findings establish that integrating multimodal deep representations with explainable AI offers a scalable, robust, and trustworthy blueprint for real-world enterprise cyber defense.
\end{spacing}
